% Existing implementation detail relocated for C7; provenance is in research/C7-RELOCATION-MAP.md.
\subsection{Software baseline}\label{app:software-baseline}

The implementation description refers to Lattice 1.7.0, released on 11 November 2026 as tag \texttt{v1.7.0} (commit \texttt{87f7bb9}), published to npm as \linebreak \code{@full-self-browsing/lattice} and \code{@full-self-browsing/lattice-cli}, and archived on Zenodo \citep{lattice}. The description was checked against commit \texttt{73a47cf} on 5 November 2026 rather than through a new measurement run; between that commit and the release, the packages changed only in error-message detail and OpenAI base-URL handling, with accompanying tests, documentation and version metadata, and their dependencies are unchanged. Both packages declare Node.js $\geq 24$; the workspace declares pnpm 10.33.1. These are manifest requirements, not the environments of every experiment. Dependency versions are recorded in \code{package.json}, the package manifests, \code{pnpm-workspace.yaml} and \code{pnpm-lock.yaml} at the release commit. A package version alone does not identify the source revision tested: revisions after the 1.6.0 release whose manifests still read 1.6.0 are not the published 1.6.0 packages.

For a source checkout, the repository documents \code{pnpm install} and \code{pnpm build}. Reproduction of an experiment requires its measured revision and retained inputs, separately from the documented baseline. An isolated checkout avoids replacing the build used by a frozen evaluation. Environment fields absent from a cohort's saved metadata are unrecorded; they cannot be inferred from the host used to build this manuscript.

\subsection{Receipt and execution details}\label{app:receipt-details}

\paragraph{DP1 Configured automatic signing.} A configured signer attempts a receipt on terminal results of \code{ai.run} and public \code{ai.runAgent}. Mint failures are flagged or thrown under the configured policy. Agent intents can override the runtime signer; terminal receipts need no tracer. Per-iteration checkpoints remain separate. Crews return a crew-start receipt, completion receipts for the root and successful child runs, and a crew-close receipt. Receipts use RFC 8785 canonicalization \citep{rfc8785}, DSSE \citep{dsse}, and Ed25519 \citep{rfc8032,bernstein2012}. Key sets track active, retired, and revoked keys; version floors reject unsupported older schemas.

\paragraph{DP2 Retain and bind content.} Receipts contain SHA-256 hashes of input artifacts and outputs and a Merkle root over artifact lineage \citep{merkle1988}. Optional v1.4 fields add task and policy hashes and an ordered list of provider-request hashes. Content is kept separately. Monetary cost is serialized as an I-JSON string \citep{rfc7493}.

\paragraph{DP3 Routing and control outcomes.} The router applies hard filters (modality, context window, structured output) and policy filters (allow and deny lists, privacy class, logging, latency, budget), then ranks survivors by a score combining cost, latency, and context headroom. Ties break on provider and model identifiers, not on catalog position. Output invariants (\code{mustCite}, \code{noPII}, \code{fieldFromTable}, \code{matches}) are evaluated by a pure function; violations are terminal and excluded from fallback.

\paragraph{DP4 Minimize disclosure; verify before loading.} Canonicalization is called only on the output of redaction. Materialization for offline replay verifies the receipt before any artifact is loaded. Input and output hashes are unsalted SHA-256.

\paragraph{Receipt contents.} A signed receipt body holds the run identifier, issue time, signing key identifier, requested and observed model, selected route, token usage and cost, contract verdict and hash, input-artifact hashes, an output hash, lineage root, and redaction manifest. For \code{ai.run} and public \code{ai.runAgent}, the output hash is set only when the run succeeds; a failure receipt commits to the run's inputs and context but not to an output. It can also carry redacted tripwire evidence or route rejection reasons. Audience, nonce, expiry, parent, and closure fields remain optional. Schema v1.4 adds optional \code{taskHash}, \code{policyHash}, and \code{providerRequestHashes}. Current signed \code{ai.run} and public \code{ai.runAgent} results populate the task and effective-policy hashes using SHA-256 over canonical JSON. The ordered request list is included only when every executed adapter supplies the submission callback and at least one body was prepared; it includes prepared attempts that later fail. Each entry hashes the exact serialized UTF-8 body handed to configured fetch, after packaging. Transport headers and the endpoint URL are excluded from that hash, and raw content is not placed in the receipt. A custom fetch can subsequently modify those bytes, and the digest proves neither remote receipt nor model consumption. Legacy receipts and adapters without complete callback coverage remain visibly unbound. Verifiers accept optional fields when absent unless the caller supplies corresponding expected digests; a version number alone does not require their presence. Routing candidates, scores, and catalog version remain outside the signed record, so DR2 and DR3 require separately retained evidence beyond these commitments.

\paragraph{Agent runs.} Public \code{ai.runAgent} now returns a signed terminal receipt on successful and failed result variants when an intent or runtime signer is configured, independently of tracing. It commits to the original task, effective policy, terminal verdict, successful output when available, and the complete observed provider-submission sequence when supported. Its default signer-failure behavior returns \code{receiptError}; the configured throw mode raises an error. Per-iteration checkpoints remain a separate tracer facility. Crew execution retains its existing receipt accounting: a crew-start root, a completion receipt for the root and each successful child run, and a crew-close receipt. Thus $N$ agents that each run once with successful children yield $N+2$ crew receipts; a failed child has no crew completion receipt. The crew-start hash covers the nested identifiers, not the whole specification, and completion receipts hash outputs. The single-agent extension does not retroactively add full task, tool, or policy commitments to these crew records. Current \code{verifyChain} enforces the closing member list; historical 1.6.x verifiers do not. \code{requireClosure} rejects missing closure, while \code{partial} permits only listed subsets with a closer. Crew work time ends before root completion and closure signing, but includes start and child signing. Failed mints are reported under the signer-error policy; a verifiable closing list names only receipts actually minted, so callers must also inspect \code{receiptError}.

\subsection{Integration interfaces}\label{app:integration-details}

\paragraph{Provider coverage.} Seven adapters (OpenAI, generic OpenAI-compatible endpoints, Anthropic, Google Gemini, xAI, OpenRouter, and LM Studio) share a provider contract. Current ordinary and streaming builders deliver supported image values as native image parts, preserve JSON values, and reject unsupported or unsatisfied transports before submission. Binary image views retain their byte ranges; URL delivery remains subject to \code{noPublicUrl}. The earlier vision probe in Section~\ref{sec:realworld-broad} records the old builder, which dropped binary and object values or sent image data as text; it is historical evidence of the defect, not a measurement of the correction. Optional \code{generation.temperature} and \code{generation.maxOutputTokens} now flow through calls, agents, and crews into provider-specific wire fields. Native OpenAI uses \code{max\_completion\_tokens}, compatible gateways use \code{max\_tokens}, and native Anthropic and Gemini use their own fields. Omitted settings preserve earlier defaults. The request digest records these serialized choices, but does not validate provider capability claims or prove that a model used an input. The router's capability catalog remains a point-in-time registry supplemented by maintained profiles.

\paragraph{External executors.} An organization that already uses another orchestration framework can call an audit helper that mints the same receipt family from externally produced request, response, and output data. This offers an incremental integration path relevant to DR5, without measuring adoption effort, and requires trusting the external executor's account of the provider call.

\paragraph{Command-line tool.} The tool verifies receipts and receipt chains, prints structured advisories as warnings or information lines, checks stored artifacts and outputs offline against the signed hashes, compares receipts structurally, and gates regressions in cost and output consistency across a corpus of receipts in continuous integration. The library can also export run events through OpenTelemetry \citep{otel}, and record Model Context Protocol tool results as artifacts \citep{mcp}.

\subsection{Measured environments and historical qualifications}\label{app:measured-environments}

The narrative quotes the harness through macros generated by \code{eval/scripts/06-report.py}. The latency and size results (Tables~\ref{tab:overhead}, \ref{tab:size}, and~\ref{tab:chain}; Figures~\ref{fig:overhead} and~\ref{fig:chain}) come from the original measurement run on Node \EvalNode{}, \EvalCpu{} (\EvalCores{} cores, \EvalMem{}~GB), platform \EvalPlatform{}, Lattice \EvalLatticeVersion{}, commit \texttt{\EvalCommit}, at \EvalDate{}. The functional experiments (tamper, replay, policy, and verifier checks) do not depend on timing. They ran against the library source at commit \texttt{\EvalFuncCommit} on Node \EvalFuncNode{}, \EvalFuncCpu{} (\EvalFuncCores{} cores, \EvalFuncMem{}~GB), platform \EvalFuncPlatform{}, at \EvalFuncDate{}, and the test counts below come from the same environment. The library tests passed \TestsLatticePassed{} of \TestsLatticeTotal{} (\TestsLatticeFailed{} failed). The command-line tests passed \TestsCliPassed{} of \TestsCliTotal{} (\TestsCliFailed{} failed). Those counts document that the suite ran. They are not the argument for the design.

The original controlled suite used mocked provider calls and read no API credential from the environment. Separate live, paid and free cohorts retain their own revisions and coverage in Appendix~\ref{app:hosted-details}; the synthetic refund experiment remains separate in Appendix~\ref{sec:synthetic-refund-case}. Dataset provenance and access limits remain with those reports and the public-record case (Section~\ref{sec:refund-case}). Underlying Ai4Privacy records and derived annotations are not redistributed.

\subsubsection{Detailed qualifications by revision}\label{app:implementation-qualifications}
The following qualifications preserve the recorded reruns, exclusions and implementation-specific boundaries. Later code does not inherit historical results.

\begin{itemize}
\item \emph{Historical mocked and free-hosted samples.} The original controlled tables use a simulated provider. Sections~\ref{sec:realworld} and~\ref{sec:realworld-broad} record a local model and free hosted models. That hosted sample is a free tier: it rate-limits, it does not represent paid latency, and one catalogued model rejects ordinary chat calls. A rerun at commit \texttt{c5f79b0} on the author's machine, started 2026-10-26 14:13 CDT, measured median Lattice processing of 1.7 to 3.0~ms (0.07 to 0.27\% of end-to-end; Liquid's p95 was 0.72\%), which is higher than the previous free-tier run. Two crews stopped at the wall-clock budget after free-tier waits and are not successes. The single-agent stage was not included in that rerun. The adapter at that historical revision sent no temperature or token cap, so its fetch wrapper added both. The current API forwards explicit generation settings before hashing the submitted body; the old timing results do not measure that path. Provider-time differences between a packed prompt and a one-line prompt are too noisy to attribute to the extra tokens.
\item \emph{Optional bindings and retained evidence.} Current v1.4 receipts support task, policy, and ordered request digests; legacy receipts and unsupported adapters can lack them. A verifier must supply expected commitments when their absence or mismatch should fail. Raw task, policy, prompt, routing candidates, and catalog evidence are not reconstructed from the receipt, and routing rationale remains outside its commitments.
\item \emph{Terminal evidence is not a complete event log.} Public \code{ai.runAgent} now returns a terminal receipt without requiring a tracer. Per-iteration checkpoints remain separate, and crew receipts retain their existing tree/accounting format. A terminal digest does not establish that every tool effect was independently witnessed or durably recorded before execution.
\item \emph{Submission is not consumption.} Current builders send supported images natively, preserve JSON, and fail unsupported transports. The earlier vision probe in Section~\ref{sec:realworld-broad} misread a red image sent as text; it was not repeated at commit \texttt{c5f79b0} and remains evidence about that earlier builder. Correct byte packaging alone does not establish image comprehension or remote consumption. The callback boundary excludes subsequent custom-fetch mutation, transport headers, and receiver attestation.
\item \emph{Unsalted hashes.} Artifact, output, task, policy, and request commitments use unsalted SHA-256. A holder who can guess the committed content can confirm it, including low-entropy decisions or inputs; these digests are not encryption.
\item \emph{Freshness binding is opt-in.} Audience, nonce, and a raised version floor are not required for verification. A receipt with no \code{expiresAt} does not expire unless the caller passes \code{maxAgeMs}. There is no replay cache. Advisories report unchecked bindings as warnings or information. They do not reject.
\item \emph{Open chains can still be tail-truncated.} Without \code{expectedCount} or a closing member list, truncation produces a warning, not a failure. Crew chains carry a closing member list, but verifiers in the 1.6.x line do not enforce it.
\item \emph{Signer failure is visible, not fail-closed by default.} If the signer throws, the run still returns its result, which carries \code{receiptError}, unless \code{onSignerError} is set to throw. A crew-start signing failure always throws, because no agent work has run. A run with no signer emits no receipt.
\item \emph{Pattern-based personal-data checks.} Historical detector measurements covered email, United States SSNs, payment cards, and phone patterns. Labels, type cues, and long numeric identifiers produced the documented misses and false positives. The public-data rows in Section~\ref{sec:realworld} were remeasured at commit \texttt{ae99fb5}; their SSN/phone trade-offs and card false positives are retained as historical findings. Current all-match and optional contextual handling require separate held-out scoring. Neither blocking a sample nor matching a marker establishes complete privacy protection or correct entity-type classification.
\item \emph{must-cite is syntactic.} By default it reads only structured citation lists. Prose checks are opt-in. It does not decide whether a cited artifact supports the prose. On ASQA dev, that default accepts a gold title in 0 of 948 answers. With \texttt{inline:~true} it accepts 4. An earlier count of 7 used a default that read inline markers.
\item \emph{Self-attestation and synthetic data.} Receipts are signed by the issuing runtime. The policy evaluation used synthetic data, and its labeled cases were chosen with the rules they test in view.
\item \emph{No human evaluation yet} (Sect.~\ref{sec:eval-study}).
\end{itemize}

\subsection{Reproduction commands}\label{app:reproduction-commands}

\paragraph{Presentation and build only.} From the repository root, the following commands check the current presentation against saved evidence without inference or new experimental measurements:
\begin{quote}\small
\code{node paper/scripts/c4-tables.mjs -{}-check}\\
\code{node paper/scripts/public-refund-tables.mjs -{}-check}
\end{quote}
Omitting \code{-{}-check} regenerates those presentation files from saved results. The synthetic presentation command, \code{node paper/scripts/refund-tables.mjs -{}-check}, additionally verifies the frozen build and bundles and repeats its saved local integrity controls; it requires the matching build and does not call a model. If its build hashes differ, use the measured revision in an isolated checkout rather than changing the frozen hashes. Build the paper with \code{make -C paper tectonic}; building the PDF does not regenerate evaluation data.

The original numerical macros and tables are generated by \code{eval/scripts/06-report.py} from saved \code{eval/results} data. The local/hosted comparison fragments in Sections~\ref{sec:realworld} and~\ref{sec:realworld-broad} are generated from \code{eval/realworld/results} by \code{eval/realworld/scripts/05-charts.py} and \code{12-charts.py}. The reporting commands, input locations and generated-file inventory are retained in \code{paper/README.md}; generated numerical fragments must not be edited by hand.

\paragraph{Experimental reruns.} The commands \code{eval/run-all.sh} and \code{pnpm bench:receipts} re-measure latency on the execution host and may replace measurements. The repository benchmark is separate from the harness tamper suite and writes \code{paper/bench/results.json}; its rendering step also updates the IEEE manuscript's benchmark blocks. These are experimental reruns, not presentation refreshes. Historical cohorts, failed cases and exclusions remain tied to their recorded revisions.
